% Joint-paper results: reward construction and low-success learning.
\subsection{Phase-Aware Feedback on LIBERO}\label{sec:libero_results}
Our LIBERO reward comparison appears in Table~\ref{tab:libero}. Primary phase-aware dense feedback increases validation success from 61.0\% to 65.3\%, compared with 61.7\% for terminal-only reward: a 3.6-percentage-point endpoint difference. Filtered dense reaches 64.7\%. Phase-only shaping reaches 62.3\%, whereas weaker-terminal dense remains at 61.0\%. The primary dense and phase-only rows differ in terminal coefficient and shaping weights; their 3.0-point gap does not isolate within-phase progress. Likewise, the table does not isolate terminal scaling. The observed ordering supports structured intermediate feedback under the tested settings, rather than a general claim that denser reward is better.
\begin{table}[t]
\caption{LIBERO-Spatial reward comparison. All configurations start at 61.0\%; $\Delta$ is in percentage points. Evaluation counts, seeds, and uncertainty remain to be completed.}
\label{tab:libero}\centering\small
\begin{tabular}{@{}lrr@{}}\toprule
Configuration & Success (\%) & $\Delta$\\\midrule
Primary dense, $\lambda=5$ & 65.3 & +4.3\\
Filtered dense $[0.1,0.9]$ & 64.7 & +3.7\\
Phase-only, $\lambda=1$ & 62.3 & +1.3\\
Terminal-only, $\lambda=5$ & 61.7 & +0.7\\
Dense, $\lambda=1$ & 61.0 & 0.0\\\bottomrule
\end{tabular}
\end{table}


In our low-success experiment, terminal-only and dense training reach 11\% and 18\%, respectively, from the same 10\% initialization (Table~\ref{tab:libero_low}). Dense training also assembles accepted batches faster and completes more optimizer updates during the observed interval. \pending{Quantify this training-log observation using timestamps, attempt counts, and matched budget endpoints.} It does not yet establish a numerical speedup or improvement per optimizer update.
\begin{table}[t]
\caption{Low-success LIBERO experiment. Both conditions use group-accuracy filtering $[0.1,1.0]$. Values are percentages; evaluation timestamps and denominators remain to be completed.}
\label{tab:libero_low}\centering\small
\begin{tabular}{@{}lrr@{}}\toprule
Evaluation & Terminal-only & Dense\\\midrule
Initial & 10 & 10\\
Intermediate & 8 & 12\\
Final & 11 & 18\\\bottomrule
\end{tabular}
\end{table}


These results support our use of intermediate feedback after basic skills exist, as in CDPR task composition. They do not determine whether dense shaping or milestone credit is preferable on the CDPR, nor establish the benefit of binary refinement or self-imitation. Those comparisons test other parts of the adaptation schedule.
